% =============================================================================
% Commandes de mise en forme du cahier des charges.
% Les fichiers de src/sections/ n'utilisent QUE ces commandes : aucune mise en
% forme écrite en dur dans les sections.
% =============================================================================

% --- Propriétés du PDF -------------------------------------------------------
\AtBeginDocument{\hypersetup{pdftitle={\docshorttitle}, pdfauthor={\docauthor}}}

% --- Titres --------------------------------------------------------------------
% Sobres : texte noir, gras, aucune icône, aucun filet.

% \cdcsection{Titre} : section numérotée « 1. Titre », gras 14 pt.
\titleformat{\section}[hang]{\fontsize{14}{17}\selectfont\bfseries}{\thesection.}{0.5em}{}
\titleformat{name=\section,numberless}[hang]{\fontsize{14}{17}\selectfont\bfseries}{}{0pt}{}
\titlespacing*{\section}{0pt}{18pt plus 4pt minus 2pt}{8pt}
\newcommand{\cdcsection}[1]{\section{#1}}

% \cdcsectionnonnum{Titre} : même style, sans numéro, présente dans la table
% des matières (historique, glossaire, annexes).
\newcommand{\cdcsectionnonnum}[1]{%
  \phantomsection
  \section*{#1}%
  \addcontentsline{toc}{section}{#1}}

% \cdcpartie{Titre} : regroupe plusieurs sections, numérotée en chiffres
% romains (« PARTIE I – TITRE »), en capitales grasses dans le texte comme
% dans la table des matières. Chaque partie commence sur une nouvelle page.
\newcounter{cdcpartie}
\newcommand{\CDCpartielibelle}[1]{Partie~\Roman{cdcpartie}~– #1}
\newcommand{\cdcpartie}[1]{%
  \stepcounter{cdcpartie}%
  \clearpage\phantomsection
  \addcontentsline{toc}{part}{\texorpdfstring{\MakeUppercase{\CDCpartielibelle{#1}}}{\CDCpartielibelle{#1}}}%
  \noindent{\fontsize{12}{15}\selectfont\bfseries\MakeUppercase{\CDCpartielibelle{#1}}}\par\nopagebreak}

% \cdcsubsection{Titre} : « 1.1 Titre », gras noir 12 pt.
\titleformat{\subsection}{\fontsize{12}{14.5}\selectfont\bfseries}{\thesubsection}{0.6em}{}
\titlespacing*{\subsection}{0pt}{12pt plus 3pt minus 2pt}{4pt}
\newcommand{\cdcsubsection}[1]{\subsection{#1}}

% \cdcannexe{Titre} : « Annexe A – Titre », à placer après \cdcsectionnonnum{Annexes}.
% Style d'une sous-section ; dans la table des matières, aligné sur les sections.
\newcounter{cdcannexe}
\newcommand{\cdcannexe}[1]{%
  \stepcounter{cdcannexe}%
  \phantomsection
  \subsection*{Annexe~\Alph{cdcannexe}~– #1}%
  \addcontentsline{toc}{subsection}{Annexe~\Alph{cdcannexe}~– #1}}

% --- Page de garde --------------------------------------------------------------
% PDF pleine page (\docgarde, défini dans meta.tex), sans en-tête ni numéro.
% Numérotation « Alph » sur la garde, puis arabe à partir de 1 pour la table des
% matières (évite les ancres hyperref en double).
\newcommand{\cdcpagedegarde}{%
  \pagenumbering{Alph}%
  \includepdf[pages=1]{\docgarde}%
  \pagenumbering{arabic}}

% --- Table des matières -----------------------------------------------------------
% Titre centré en gras ; parties en capitales grasses sur un filet fin ; sections en
% gras avec points de conduite ; sous-sections en romain, en retrait.
% Aucune icône.
\fancypagestyle{tdm}{%
  \fancyhf{}%
  \fancyfoot[C]{\small Page~\thepage{} sur~\pageref*{LastPage}}%
  \renewcommand{\headrulewidth}{0pt}%
  \renewcommand{\footrulewidth}{0pt}}

\contentsmargin{2.2em}
\titlecontents{part}[0pt]
  {\addvspace{10pt}\bfseries}
  {}{}{}
  [\endgraf\vspace{2pt}{\color{rulegray}\hrule height 0.4pt}\vspace{1pt}]
\titlecontents{section}[1.8em]
  {\addvspace{3pt}\bfseries}
  {\contentslabel{1.8em}}
  {\hspace*{-1.8em}}
  {\mdseries\titlerule*[0.6pc]{.}\bfseries\contentspage}
\titlecontents{subsection}[4.3em]
  {\addvspace{1pt}}
  {\contentslabel{2.5em}}
  {\hspace*{-2.5em}}
  {\titlerule*[0.6pc]{.}\contentspage}

\newcommand{\cdctabledesmatieres}{%
  \clearpage
  \pagestyle{tdm}%
  \phantomsection\pdfbookmark[1]{Table des matières}{tdm}%
  {\centering\fontsize{18}{22}\selectfont\bfseries Table des matières\par}%
  \vspace{10pt}%
  {\setlength{\parskip}{0pt}\csname @starttoc\endcsname{toc}}%
  \clearpage
  \pagestyle{cdc}}

% --- En-tête de tableau (marine, gras) ------------------------------------------
\newcommand{\cdcth}[1]{\textcolor{primary}{\textbf{#1}}}

% --- Priorités MoSCoW (distinguées par le texte et la graisse, pas la couleur) --
\newcommand{\prioMust}{\textbf{Must have}}
\newcommand{\prioShould}{\textbf{Should have}}
\newcommand{\prioCould}{Could have}
\newcommand{\prioWont}{\textit{Won’t have}}

% --- Fiche descriptive d'une fonction (grille du professeur) ---------------------
% \fichefonction{ID}{Intitulé}{Objectif}{Description}{Contraintes / règles de gestion}{Priorité}
% Tableau insécable : en-tête « Fonction ID – Intitulé » en marine, puis les
% quatre lignes de la grille (libellés en gras dans la colonne de gauche).
\newlength{\cdccollibelle}\setlength{\cdccollibelle}{3.6cm}
\newcommand{\fichefonction}[6]{%
  \par\addvspace{8pt}\noindent
  \begin{tabularx}{\linewidth}{@{}>{\raggedright\arraybackslash\bfseries}p{\cdccollibelle}%
                                  >{\raggedright\arraybackslash}X@{}}
    \toprule
    \multicolumn{2}{@{}l@{}}{\cdcth{Fonction #1~– #2}} \\
    \midrule
    Objectif & #3 \\ \addlinespace[3pt]
    Description & #4 \\ \addlinespace[3pt]
    Contraintes / règles de gestion & #5 \\ \addlinespace[3pt]
    Niveau de priorité & #6 \\
    \bottomrule
  \end{tabularx}\par\addvspace{8pt}}

% --- Tableau d'exigences (longtable, en-tête répété) ---------------------------
% \begin{tableexigences}
%   EF-01 & Intitulé & \prioMust & Source \\
% \end{tableexigences}
\newlength{\cdccolid}\setlength{\cdccolid}{1.7cm}
\newlength{\cdccolprio}\setlength{\cdccolprio}{2.4cm}
\newlength{\cdccolsource}\setlength{\cdccolsource}{3cm}
\newlength{\cdccolexig}

\newenvironment{tableexigences}{%
  \setlength{\cdccolexig}{\dimexpr\linewidth-\cdccolid-\cdccolprio-\cdccolsource-6\tabcolsep\relax}%
  \renewcommand{\arraystretch}{1.15}%
  \setlength{\LTpre}{6pt}\setlength{\LTpost}{6pt}%
  \begin{longtable}{@{}>{\raggedright\arraybackslash}p{\cdccolid}%
                      >{\raggedright\arraybackslash}p{\cdccolexig}%
                      >{\raggedright\arraybackslash}p{\cdccolprio}%
                      >{\raggedright\arraybackslash}p{\cdccolsource}@{}}
    \toprule
    \cdcth{ID} & \cdcth{Exigence} & \cdcth{Priorité} & \cdcth{Source} \\
    \midrule
    \endfirsthead
    \toprule
    \cdcth{ID} & \cdcth{Exigence} & \cdcth{Priorité} & \cdcth{Source} \\
    \midrule
    \endhead
    \midrule
    \multicolumn{4}{r@{}}{\small\textit{Suite page suivante}} \\
    \endfoot
    \bottomrule
    \endlastfoot
}{%
  \end{longtable}}

% --- Tableau générique (booktabs, en-tête marine répété, sécable) ----------------
% \begin{cdctableau}{colonnes}{\cdcth{A} & \cdcth{B}}
%   cellule & cellule \\
% \end{cdctableau}
% Colonnes : P{largeur} = largeur fixe ; L = largeur libre, les colonnes L se
% partagent la place restante (calculée ci-dessous). Uniquement P et L.
% (xltabular est évité : avec booktabs, il provoque « Infinite glue shrinkage »
% aux coupures de page.)
\newlength{\cdcreste}
\newlength{\CDCfixe}
\newcount\CDCnL
\newcount\CDCncol
\makeatletter
\def\CDCstop{\CDCstop}
\def\CDCscan#1{%
  \ifx#1\CDCstop\expandafter\@gobble\else\expandafter\@firstofone\fi
  {\advance\CDCncol\@ne
   \ifx#1P\expandafter\CDCscanP\else\advance\CDCnL\@ne\expandafter\CDCscan\fi}}
\def\CDCscanP#1{\advance\CDCfixe#1\relax\CDCscan}
\newcommand{\CDCcalculreste}[1]{%
  \setlength{\CDCfixe}{0pt}\CDCnL=0 \CDCncol=0
  \CDCscan#1\CDCstop
  \setlength{\cdcreste}{\dimexpr\linewidth-\CDCfixe-2\tabcolsep*(\CDCncol-1)\relax}%
  \ifnum\CDCnL>1 \divide\cdcreste\CDCnL\fi}
\makeatother
\newcolumntype{L}{>{\raggedright\arraybackslash}p{\cdcreste}}
\newcolumntype{P}[1]{>{\raggedright\arraybackslash}p{#1}}
\newenvironment{cdctableau}[2]{%
  \CDCcalculreste{#1}%
  \setlength{\LTpre}{6pt}\setlength{\LTpost}{6pt}%
  \begin{longtable}{@{}#1@{}}
    \toprule #2 \\ \midrule \endfirsthead
    \toprule #2 \\ \midrule \endhead
    \midrule \endfoot
    \bottomrule \endlastfoot}%
  {\end{longtable}}

% \cdcsep : espace entre deux rangées d'un tableau (en début de rangée).
\newcommand{\cdcsep}{\addlinespace[4pt]}
% \groupeexigences{Titre} : rangée d'intertitre dans tableexigences.
\newcommand{\groupeexigences}[1]{\addlinespace[4pt]\multicolumn{4}{@{}l}{\textbf{#1}}}

% Liste à puces à l'intérieur d'une cellule de tableau (sans espace parasite).
% La minipage [t] aligne la première puce sur la ligne de base de la rangée.
\newenvironment{cdccellitems}%
  {\begin{minipage}[t]{\linewidth}\raggedright
   \begin{itemize}[label=\textbullet, leftmargin=1em, itemsep=0pt, parsep=0pt,
     topsep=0pt, partopsep=0pt]}%
  {\end{itemize}\end{minipage}}

% --- Étapes numérotées (scénarios) ----------------------------------------------
\newenvironment{cdcetapes}%
  {\begin{enumerate}[label=\arabic*., leftmargin=1.8em, itemsep=3pt, parsep=0pt, topsep=2pt]}%
  {\end{enumerate}}

% --- Maquettes (dessins TikZ placés dans assets/maquettes/) ------------------------
% Filaires, en noir et gris uniquement, texte en 10 pt minimum.
% \cdcmaquette{assets/maquettes/fichier}{Légende}            : une maquette centrée
% \cdcmaquettes{fichier1}{Légende 1}{fichier2}{Légende 2}     : deux côte à côte
\tikzset{
  mq ecran/.style={draw=black, line width=0.8pt, rounded corners=8pt},
  mq barre/.style={fill=black!8, draw=none},
  mq texte/.style={anchor=north west, inner sep=0pt, font=\small, align=left},
  mq titre/.style={mq texte, font=\small\bfseries},
  mq gris/.style={mq texte, text=black!60},
  mq champ/.style={draw=black!45, anchor=north west, minimum height=0.62cm,
    inner xsep=5pt, font=\small, text=black!60, align=left, text width=5.55cm},
  mq bouton/.style={fill=black, text=white, font=\small\bfseries,
    rounded corners=3pt, minimum height=0.7cm, inner xsep=6pt},
  mq bouton2/.style={draw=black, text=black, font=\small\bfseries,
    rounded corners=3pt, minimum height=0.7cm, inner xsep=6pt},
  mq cadre/.style={draw=black!45, rounded corners=3pt},
  mq photo/.style={draw=black!45, fill=black!5},
}
% \mqphoto{x}{y}{taille} : vignette d'image (rectangle barré).
\newcommand{\mqphoto}[3]{%
  \draw[mq photo] (#1,#2) rectangle ++(#3,-#3);
  \draw[black!30] (#1,#2) -- ++(#3,-#3) (#1,#2-#3) -- ++(#3,#3);}

\newcounter{cdcmaquette}
\newcommand{\CDClegendemaquette}[1]{%
  \stepcounter{cdcmaquette}%
  {\small\textbf{Maquette~\arabic{cdcmaquette}}~– #1}}
\newcommand{\cdcmaquette}[2]{%
  \par\addvspace{10pt}%
  \begin{center}
    \input{#1}\par\vspace{6pt}%
    \CDClegendemaquette{#2}
  \end{center}
  \par\addvspace{10pt}}
\newcommand{\cdcmaquettes}[4]{%
  \par\addvspace{10pt}\noindent
  \begin{minipage}[t]{0.49\linewidth}\centering
    \input{#1}\par\vspace{6pt}\CDClegendemaquette{#2}
  \end{minipage}\hfill
  \begin{minipage}[t]{0.49\linewidth}\centering
    \input{#3}\par\vspace{6pt}\CDClegendemaquette{#4}
  \end{minipage}
  \par\addvspace{10pt}}

% --- Encadrés note et attention -------------------------------------------------
% Fond blanc, filet marine à gauche, sans icône.
\tcbset{cdcencadre/.style={enhanced, breakable, colback=white, colframe=primary,
  boxrule=0pt, leftrule=2pt, arc=0pt, outer arc=0pt,
  left=8pt, right=0pt, top=3pt, bottom=3pt, before skip=8pt, after skip=8pt}}

\newtcolorbox{note}{cdcencadre,
  before upper={\textbf{Note}~:\ }}
\newtcolorbox{attention}{cdcencadre,
  before upper={\textbf{Attention}~:\ }}

% --- Listes compactes -----------------------------------------------------------
\newenvironment{cdcitems}%
  {\begin{itemize}[label=\textbullet, leftmargin=1.1em, itemsep=3pt, parsep=0pt, topsep=2pt]}%
  {\end{itemize}}
% Sous-liste à tirets, à placer dans un \item de cdcitems ou de cdcetapes.
\newenvironment{cdcsousitems}%
  {\begin{itemize}[label=–, leftmargin=1.1em, itemsep=0pt, parsep=0pt, topsep=1pt]}%
  {\end{itemize}}

% --- Liste de définitions (glossaire) ----------------------------------------------
% \begin{cdcdefinitions}[largeur des libellés] \item[Terme] Définition. \end{cdcdefinitions}
% Largeur par défaut : 4.6cm (glossaire).
\newenvironment{cdcdefinitions}[1][4.6cm]%
  {\begin{description}[style=multiline, leftmargin=#1, labelwidth=\dimexpr#1-0.2cm\relax,
     font=\normalfont\bfseries, itemsep=2pt, topsep=2pt]}%
  {\end{description}}

% --- Intertitre de paragraphe ---------------------------------------------------
% \cdcintertitre{Titre} : gras, non numéroté, absent de la table des matières.
\newcommand{\cdcintertitre}[1]{\par\addvspace{6pt}\noindent\textbf{#1}\par\nopagebreak}

% --- Ligne « Libellé : texte » ---------------------------------------------------
\newcommand{\cdcdetail}[2]{\par\noindent\textbf{#1}~: #2\par}

% --- Historique des versions ------------------------------------------------------
% \historique{
%   \versionligne{0.1}{15/10/2026}{Auteur}{Création du document}
% }
\newcommand{\versionligne}[4]{#1 & #2 & #3 & #4 \\}
\newcommand{\historique}[1]{%
  \par\addvspace{4pt}\noindent
  \begin{tabularx}{\linewidth}{@{}l l l >{\raggedright\arraybackslash}X@{}}
    \toprule
    \cdcth{Version} & \cdcth{Date} & \cdcth{Auteur} & \cdcth{Modifications} \\
    \midrule
    #1
    \bottomrule
  \end{tabularx}\par\addvspace{4pt}}

% --- Liens et e-mails (bleu de lien) ---------------------------------------------
\newcommand{\cdclien}[1]{\textcolor{lien}{\url{#1}}}
\newcommand{\cdcemail}[1]{\href{mailto:#1}{\textcolor{lien}{#1}}}
